\scanpage{238}

IV -- \mspage{9} -- \hfill n°338

que $\underline m^n\subset\underline q$ (Chap.III, §2, n°2, cor. 2 du lemme 2).

3) Dans un anneau \uline{principal} $A$, les idéaux primaires sont $(0)$ et ceux de la forme $(p^n)$ avec $p$ extrémal et $n\geq1$ : ceci résulte de l'ex. 2.

4) Les puissances d'un idéal premier quelconque $\underline p$ ne sont pas en général des idéaux primaires (exerc.\quad). Il existe des idéaux primaires qui ne sont pas des puissances d'idéaux premiers (exerc.\quad).

\margin{\unsure{2}}

PROPOSITION 2.- \uline{Soient} $M$ \uline{un module sur un anneau noethérien} $A$, $\underline p$ \uline{un idéal premier de} $A$, \uline{et} $(Q_i)_{i\in I}$ \uline{une famille finie non vide de sous-modules $\underline p$-primaires de} $M$. \uline{Alors} $\bigcap_{i\in I}Q_i$ \uline{est $\underline p$-primaire}.

En effet $M/\bigcap_{i\in I}Q_i$ est isomorphe à un sous-module non nul de la somme directe $\sum_{i\in I}M/Q_i$. Or $\operatorname{Ass}(\sum_{i\in I}M/Q_i)=\bigcup_{i\in I}\operatorname{Ass}(M/Q_i)=\{\underline p\}$ (§1, cor.1 de la prop. 3). Donc $\operatorname{Ass}(M/\bigcap_{i\in I}Q_i)=\{\underline p\}$ (§1, prop.3).

PROPOSITION 3.- \uline{Soient} $A$ \uline{un anneau noethérien}, $S$ \uline{une partie multiplicative de} $A$, $\underline p$ \uline{un idéal premier de} $A$, $M$ \uline{un $A$-module}, $N$ \uline{un sous-module} $M$, \uline{et} $i=i_M^S$ \uline{l'application canonique de} $M$ \uline{dans} $S^{-1}M$.

\margin{\figure{\begin{tikzpicture}\draw (0,.05)--(.45,0);\draw (.38,.06)--(.45,0)--(.39,-.05);\end{tikzpicture}}}

a) \uline{On suppose} $\underline p\cap S\neq\varnothing$. \uline{Si} $N$ \uline{est $\underline p$-primaire}, \uline{on a} $S^{-1}N=S^{-1}M$.

b) \uline{On suppose} $\underline p\cap S=\varnothing$. \uline{Pour que} $N$ \uline{soit $\underline p$-primaire}, \uline{il faut et il suffit que} $N$ \uline{soit de la forme} $i^{-1}(N')$, \uline{où} $N'$ \uline{est un sous-module $(S^{-1}\underline p)$-primaire de} $S^{-1}M$. \uline{On a alors} $N'=S^{-1}N$.

Si $\underline p\cap S\neq\varnothing$ et si $N$ est $\underline p$-primaire, on a $\operatorname{Ass}(S^{-1}(M/N))=\varnothing$ (§1, prop.5), donc $S^{-1}M/S^{-1}N=S^{-1}(M/N)=(0)$ (§1, cor.1 de la prop.2), et par conséquent $S^{-1}N=S^{-1}M$.

Supposons $\underline p\cap S=\varnothing$. Si $N$ est $\underline p$-primaire, on a $\operatorname{Ass}(S^{-1}(M/N))=\{S^{-1}\underline p\}$ (§1, prop.5), donc le sous-module $N'=S^{-1}N$ de $S^{-1}M$ est $(S^{-1}\underline p)$-primaire ; en outre, si $s\in S$ et $m\in M$ sont tels que $sm\in N$, on a $m\in N$ (car l'homothétie définie par $s$ dans $M/N$ est injective), d'où $N=i^{-1}(N')$ (Chap.II,

\mspage{237}
